\chapter{Introduction}
\label{ch:introduction}

\section{Motivation}
\label{sec:motivation}

Mobile software is increasingly assumed to be "always connected." Modern
frameworks, backend-as-a-service platforms, and UI toolkits are built around
the request/response cycle: an action in the interface triggers a network
call, and the interface waits for -- or reacts to -- the server's answer.
This assumption is reasonable for a large class of consumer applications
used indoors, in urban areas, on stable Wi-Fi or 4G/5G connections. It is a
poor assumption for the growing class of applications used \emph{in the
field}: on construction sites, in warehouses and basements, in rural areas,
during disaster response, in agriculture, and -- the scenario adopted in
this thesis -- by facility-maintenance staff and building occupants moving
through a large university campus.

In all of these environments, connectivity is not binary. It is
intermittent (available in some rooms, absent in others), variable in
quality (a slow, high-latency link that behaves worse than no link at all
for a synchronous request), and sometimes simply unavailable for extended
periods (elevators, basements, remote buildings, network outages). An
application that treats the network as always present will, in these
conditions, either block the user interface while waiting for a server that
cannot be reached, silently drop user input, or crash outright. None of
these outcomes are acceptable for an application whose entire purpose is to
let a user \emph{capture something at the moment it happens} -- a broken
projector noticed while walking past a lecture hall, a damaged chair found
during a room inspection -- because if the app does not let the user record
the observation there and then, the observation is very often never
recorded at all.

\emph{Offline-first} is the name given to a family of design and
engineering practices that respond to this reality by inverting the default
assumption: the local device, not the remote server, is treated as the
primary source of truth for user-generated data, and the network is treated
as an optional, asynchronous channel used to propagate that data to a
backend and to other devices. An offline-first application is, by
definition, fully usable with no network connection; when connectivity is
available, a synchronization process reconciles local and remote state in
the background, ideally without the user having to do anything beyond being
informed of the outcome.

Building an offline-first application well is a genuinely different
engineering problem from building a "normal" client-server application with
some caching bolted on. It requires an explicit local persistence model,
an explicit representation of "this record has not yet reached the
server," an explicit strategy for deciding \emph{when} to attempt
synchronization, and an explicit -- even if intentionally simple --
strategy for what happens when the same record has been modified in two
places before either side has seen the other's change. Consumer
applications like Google Docs, Notion, or messaging apps solve versions of
this problem with considerable engineering investment (operational
transformation, CRDTs, vector clocks). The central question this thesis
asks is more modest and more practical: \emph{what is the smallest,
most maintainable offline-first architecture that is sufficient for a
CRUD-style field-reporting application}, and what engineering trade-offs
does that architecture force the developer to make explicit?

\section{Offline-First Applications}
\label{sec:offline-first-background}

The term "offline first" was popularized in mobile and web development
around 2013--2014 as a reaction to the "mobile first, then make it work
offline" pattern, which tends to produce applications where offline support
is a thin, unreliable layer retrofitted onto a fundamentally online
architecture~\cite{offlinefirst2013}. The offline-first position instead
argues that offline behaviour should be a first-class design constraint
from the start of a project, because retrofitting it later requires
touching almost every part of the data layer.

At the core of offline-first engineering is a shift in where reads and
writes are directed. \Cref{fig:online-vs-offline} contrasts the two
architectures at a high level. In a conventional "online-only" design
(\cref{fig:online-vs-offline}, left), the UI issues a request directly to
the backend and the operation's success or failure is determined
synchronously by the network round trip. In an offline-first design
(\cref{fig:online-vs-offline}, right), the UI only ever talks to the local
data store; a separate synchronization component is responsible for
propagating local changes to the backend when a network path exists, and
for pulling remote changes back down. The user-visible effect of a create,
update, or delete operation is therefore always fast and always available,
because it depends only on local disk I/O, never on the network.

\begin{figure}[htbp]
  \centering
  \begin{tikzpicture}[
      node distance=1.1cm and 1.6cm,
      box/.style={draw, rounded corners, minimum width=2.6cm, minimum height=1cm, align=center, font=\small},
      arrow/.style={-Latex, thick}
    ]
    % Online-only
    \node[box, fill=blue!8] (ui1) {UI};
    \node[box, fill=orange!12, below=of ui1] (backend1) {Backend API};
    \node[box, fill=orange!6, below=of backend1] (db1) {Remote DB};
    \draw[arrow] (ui1) -- node[right, font=\scriptsize]{sync. request} (backend1);
    \draw[arrow] (backend1) -- (db1);
    \node[font=\small\bfseries, above=0.3cm of ui1] {Online-only};

    % Offline-first
    \node[box, fill=blue!8, right=4.4cm of ui1] (ui2) {UI};
    \node[box, fill=green!10, below=of ui2] (local2) {Local DB\\ (SQLite / Room)};
    \node[box, fill=violet!10, below=1.6cm of local2] (sync2) {Sync Manager};
    \node[box, fill=orange!12, right=2.2cm of sync2] (backend2) {Backend\\ (Firestore)};
    \draw[arrow] (ui2) -- node[right, font=\scriptsize]{read/write} (local2);
    \draw[arrow] (local2) -- node[right, font=\scriptsize]{queued rows} (sync2);
    \draw[arrow, <->] (sync2) -- node[above, font=\scriptsize]{async. sync} (backend2);
    \node[font=\small\bfseries, above=0.3cm of ui2] {Offline-first};
  \end{tikzpicture}
  \caption{Data flow in an online-only architecture versus an offline-first
  architecture. In the offline-first case the UI never blocks on the
  network: every read and write is served by the local database, and
  synchronization runs as an independent, asynchronous process.}
  \label{fig:online-vs-offline}
\end{figure}

This shift has consequences that ripple through the whole stack:

\begin{itemize}
  \item \textbf{Local persistence becomes a primary concern, not a cache.}
  The local database must be able to represent every state a record can be
  in, including states that have no meaning on the server (e.g. "created
  locally, not yet uploaded").
  \item \textbf{Identity must be decided before the server is involved.}
  If a report is created offline, it needs an identifier the moment it is
  created, before any server-assigned identifier exists. This typically
  means client-generated identifiers (UUIDs) coexisting with
  server-assigned identifiers.
  \item \textbf{Synchronization is a background process, not a request
  handler.} It must be triggered by connectivity changes and/or periodic
  timers, must be resumable after being interrupted, and must be
  idempotent -- running it twice must not create duplicate data.
  \item \textbf{Conflicts become possible, and must be resolved by policy,
  not by accident.} Two devices -- or a device and a server-side edit --
  can modify the same record while disconnected from each other. The
  application must define, even if only in the simplest possible way, what
  happens when those two versions are reconciled.
  \item \textbf{The user needs visibility into synchronization state.}
  Because writes are decoupled from network availability, the UI must
  communicate whether a given piece of data is safely stored on the server
  or still only local, or the offline-first property becomes a source of
  anxiety rather than resilience.
\end{itemize}

These five concerns map directly onto the "Research/engineering questions"
that motivate this thesis (\cref{sec:objectives}), and onto the chapters
that follow: \cref{ch:requirements} turns them into concrete requirements,
\cref{ch:design} turns the requirements into an architecture, and
\cref{ch:implementation} shows how that architecture is realised in code.

\section{Field Data Collection as a Case Study}
\label{sec:case-study}

Field data collection is a natural and well-studied application area for
offline-first design~\cite{brunette2013odk,hartung2010open}. It shares the
same shape across many domains -- agriculture, epidemiology, environmental
monitoring, humanitarian logistics, facility maintenance -- namely, a field
worker observes something, records structured data about it (often with a
photo and a location), and that data needs to eventually reach a central
system where it can be aggregated, triaged, and acted upon. The
observation happens where it happens; the network may or may not be
available at that exact place and time.

This thesis instantiates the field data collection problem with a scenario
that is small enough to implement and evaluate within the scope of a
master's thesis, yet representative of the harder domains listed above: a
mobile application, referred to throughout this document as
\textbf{CampusFix}, that lets students, faculty, and staff of a university
campus report facility problems as they encounter them. Representative
report categories include:

\begin{itemize}
  \item a broken or missing projector in a lecture hall;
  \item a damaged chair or desk;
  \item a Wi-Fi outage or coverage problem in a specific building/room;
  \item an air-conditioning or heating malfunction;
  \item missing or malfunctioning laboratory equipment.
\end{itemize}

A user of CampusFix is frequently \emph{exactly} in the situation
offline-first design targets: standing in a basement lab with no signal,
noticing a broken piece of equipment, and wanting to log it before moving
on and forgetting the details. If the application requires a live
connection to save that report, the report is lost, or the user is forced
to remember to re-open the app later -- which, in practice, means the
report is usually never filed. CampusFix must therefore let the user
complete the entire "create a report" flow with zero network availability,
and must reliably deliver that report to the facilities-management backend
once the device regains connectivity, without requiring the user to
remember to do anything.

\section{Problem Statement}
\label{sec:problem-statement}

Given the scenario above, this thesis addresses the following problem:
\emph{design, implement, and evaluate a mobile application that lets a user
create, view, and manage field reports with full functionality while
offline, that synchronizes those reports with a backend transparently and
reliably once connectivity returns, that communicates synchronization state
to the user at all times, and that resolves the simple conflicts that can
arise when a report is modified both locally and on the server}, using a
technology stack -- Kotlin Multiplatform for the client, SQLite/Room for
local storage, and Firebase (Firestore and Authentication) for the backend
-- that is realistic for a small engineering team to adopt and maintain.

\section{Objectives and Contributions}
\label{sec:objectives}

The thesis pursues the following concrete objectives, mirroring the
engineering questions posed in the project brief:

\begin{enumerate}[label=\textbf{O\arabic*.}]
  \item \textbf{Usability without connectivity.} Determine what UI/UX and
  architectural patterns keep a mobile app fully usable when internet
  connectivity is absent, and implement them for the report creation,
  listing, and detail flows.
  \item \textbf{Local data model.} Determine what data must be stored
  locally (and in what form) so that the application never depends on a
  network round trip for a core operation, and design a local schema that
  supports this.
  \item \textbf{Synchronization triggers.} Determine when and how
  synchronization should be triggered -- on connectivity change, on app
  start, periodically, or manually -- and implement a connectivity-aware
  synchronization module.
  \item \textbf{Sync-state visibility.} Determine how to communicate
  upload state to the user (pending / synced / failed) so that offline
  operation is legible rather than opaque.
  \item \textbf{Conflict resolution.} Determine what conflict-resolution
  strategy is acceptable for this class of application, and implement the
  simplest strategy that meets that bar.
\end{enumerate}

Correspondingly, the thesis contributes:

\begin{enumerate}[label=\textbf{C\arabic*.}]
  \item A requirements analysis (\cref{ch:requirements}) that separates
  functional requirements from the offline-specific non-functional
  requirements that are often left implicit in mobile app specifications.
  \item A layered system architecture and synchronization algorithm
  (\cref{ch:design}) for a Kotlin Multiplatform, Room-backed,
  Firestore-synchronized field-reporting application.
  \item A working implementation (\cref{ch:implementation}) of that
  architecture, shared across Android and iOS from a single Kotlin
  codebase, described down to the level of the actual data access objects,
  synchronization loop, and connectivity observers used.
  \item An evaluation (\cref{ch:evaluation}) of the system's offline
  behaviour, synchronization performance, conflict handling, and usability,
  together with an honest discussion of the current implementation's
  limitations, in particular that synchronization today is a one-directional
  push rather than a full bidirectional merge.
\end{enumerate}

\section{Thesis Structure}
\label{sec:structure}

The remainder of this thesis is organized as follows.
\Cref{ch:requirements} analyzes the functional and non-functional
requirements of CampusFix, with particular attention to the requirements
that are specific to offline operation. \Cref{ch:design} presents the
system architecture: the mobile application's internal layering, the local
database schema, the backend structure, and the design of the
synchronization strategy, including the conflict-resolution policy.
\Cref{ch:implementation} describes the implementation of that design using
Kotlin Multiplatform, Compose Multiplatform, Room, and Firebase, walking
through the local database, the connectivity monitor, the synchronization
module, and the user-facing screens. \Cref{ch:evaluation} evaluates the
resulting system against six scenarios covering offline creation, delayed
and failed synchronization, large sync queues, conflicting edits, and
end-user usability. Finally, \cref{ch:conclusions} summarizes the findings,
discusses the limitations of the current implementation, and outlines
directions for future work.
